% DRAFT — why the three-band skeleton is what it is, derived from the paper's OWN Prop. 1.
% NOT \input anywhere. Compare against the evidence-audited draft before merging.
%
% The chain, in one line:
%   Prop.1 (V_w -> span{1,Delta})  =>  the supplied-basis rank of a slot is a FUNCTION OF wL
%   =>  compressing changes wL by 4^-m, so compression MOVES a slot along the collapse curve
%   =>  the three bands are the three regimes of that curve, and only the middle one is free.

\subsection{Why the three-band family works}
\label{sec:threeband}

The constructions in \S\ref{sec:construction} are not the only allocations in use;
practically every deployed long-context recipe --- position interpolation, YaRN, and
the mixed-radix family of \citet{tian2026mrrope} --- keeps the fastest slots at their
native frequencies, leaves the slowest slots maximally compressed, and treats the
band between them as the design variable, with its edges fixed by a turn-count rule.
Proposition~\ref{prop:collapse} explains why the three bands take the roles they do.

A slot's positional content is governed by the product $\omega L$: its supplied
subspace is $\operatorname{span}\{\cos(\omega\Delta),\sin(\omega\Delta)\}$ measured
under a separation law on $[0,L]$, and as $\omega L\to0$ that subspace collapses onto
$\operatorname{span}\{1,\Delta\}$. \emph{Compression moves a slot along this curve},
because $m$ replaces $\omega$ by $\omega 4^{-m}$: the quantity that decides how much
positional information a slot supplies is $\omega L 4^{-m}$, not $\omega$ alone. The
three bands are the three regimes of that single curve.

\paragraph{The fast band is not a design variable.} Slots with $\omega L\gg1$ are far
from the collapsed limit, and their $\cos$/$\sin$ pairs supply distinguishable
directions within the window. Because compression multiplies $\omega L$ by $4^{-m}$,
moving such a slot by even one full compression pushes it toward the collapsed
regime, where its contribution to within-window resolution shrinks. The measured
consequence is the in-window cost that accompanies every allocation change in this
paper, and it is why no method in this family touches the fast band.

\paragraph{The slow band is where compression is nearly free.} Proposition~\ref{prop:collapse}
is quantitative: in the standard configuration, the $23$ pairs with $\omega L\le1$
supply $46$ coordinate dimensions but a block-whitened effective rank of
$r_2=2.00013$, i.e. a mean pairwise squared canonical overlap of $\bar c=0.99993$.
Those slots already carry almost no distinct positional information, so spending
compression on them removes redundancy rather than resolution --- which is why the
slow band can be driven to a full compression, and beyond, at little in-window cost.

\paragraph{The middle band is the only free direction.} Between those regimes lies the
band where $\omega L$ is of order one and the collapse is partial. This is the only
region where compression can be exchanged against range without either discarding
resolution (fast band) or reallocating directions that were already indistinguishable
(slow band). Its edges are therefore set by the point at which a slot stops resolving
the training window, which is what a turn-count rule estimates.

\paragraph{Consequence for the interior shape.} If the middle band is where the
collapse is gradual, then the interior profile decides \emph{how many slots sit at
which point of the curve} --- a reallocation of supplied rank across depths. That is a
real degree of freedom, and the experiments above show it changes learned behaviour at
fixed range; but the same experiments show that which profile is preferable depends on
the weights, which is why the shape cannot be ranked from the schedule alone
(Theorem~\ref{thm:non-identifiability}). This is also the reason the amplitude of the
compression is set by the target rather than fixed: extending the target lengthens $L$,
which moves the whole curve, so the band edges and the plateau move with it.
